\section{Models and methods} \label{sec:methods}

This section sets up a unified notation for dynamic discrete choice and inverse reinforcement learning, then walks through the twelve estimators that \pkg{econirl} implements. Each estimator subsection gives the objective and the variants exposed in the API. We include compact pseudocode for the four estimators whose outer structure is non-trivial, namely NFXP, TD-CCP, MCE-IRL, and AIRL. Proofs and convergence arguments are left to the cited papers.

\subsection{Unified notation}

A dynamic discrete choice problem is a tuple $(\mathcal{S}, \mathcal{A}, P, r, \beta)$. The state space $\mathcal{S}$ is finite with cardinality $|\mathcal{S}|$. The action space $\mathcal{A}$ is finite with cardinality $|\mathcal{A}|$. Transition probabilities are $P(s' \mid s, a)$ for every triple in $\mathcal{S} \times \mathcal{A} \times \mathcal{S}$. The discount factor $\beta$ lies strictly between zero and one. The flow reward is $r(s, a)$. The decision maker observes the state, chooses an action, receives flow reward, and transitions to a new state according to $P$. A type-one extreme value error with location zero and scale $\sigma$ is added to the deterministic flow utility before the action is taken, following the convention in \citet{rust1987}. The conditional value function is
\begin{equation}
v(s, a) = r(s, a) + \beta \sum_{s' \in \mathcal{S}} P(s' \mid s, a) V(s'),
\label{eq:v}
\end{equation}
where the integrated value $V(s)$ is the soft maximum of $v(s, a)$ across actions,
\begin{equation}
V(s) = \sigma \log \sum_{a \in \mathcal{A}} \exp\left(\frac{v(s, a)}{\sigma}\right) + \sigma \gamma,
\label{eq:V}
\end{equation}
with $\gamma$ the Euler-Mascheroni constant. The optimal choice probability follows the multinomial logit form
\begin{equation}
\pi(a \mid s) = \frac{\exp(v(s, a) / \sigma)}{\sum_{a' \in \mathcal{A}} \exp(v(s, a') / \sigma)}.
\label{eq:pi}
\end{equation}
The structural primitives parametrize the reward as $r(s, a) = \theta^\top \phi(s, a)$, where $\phi$ is a known feature mapping and $\theta$ is the unknown parameter vector to be estimated. The data are a panel of $N$ individuals observed for $T_i$ periods each, with realized states $s_{it}$ and actions $a_{it}$.

The objects below recur in the estimator subsections that follow. The log-likelihood of the panel under choice probabilities $\pi(\cdot \mid \cdot ; \theta)$ is
\begin{equation}
\ell(\theta) = \sum_{i=1}^N \sum_{t=1}^{T_i} \log \pi(a_{it} \mid s_{it}; \theta).
\label{eq:loglik}
\end{equation}
The soft Q-function is the conditional value function relabeled when the reward is unknown, $Q(s, a) = r(s, a) + \beta \sum_{s'} P(s' \mid s, a) V(s')$, with $V(s) = \sigma \log \sum_a \exp(Q(s, a) / \sigma) + \sigma \gamma$. The discounted state-action occupancy under policy $\pi$ is
\begin{equation}
\mu_\pi(s, a) = (1 - \beta) \sum_{t=0}^{\infty} \beta^t \Pr_\pi(s_t = s, a_t = a).
\label{eq:occ}
\end{equation}
The empirical state-action occupancy in the panel is $\hat\mu(s, a) = n^{-1} \sum_{i, t} \mathbb{1}\{s_{it} = s, a_{it} = a\}$ where $n = \sum_i T_i$. The empirical feature expectation is $\hat \E[\phi] = n^{-1} \sum_{i, t} \phi(s_{it}, a_{it})$, and the model-implied feature expectation is $\E_\pi[\phi] = \sum_{s, a} \mu_\pi(s, a) \phi(s, a)$.

The protocol implemented in \pkg{econirl} is uniform across all twelve estimators. Every estimator receives a \code{Panel} object with the realized data, a \code{UtilityFunction} object that defines $\phi$ and the parameter names, a \code{DDCProblem} object that fixes $\beta$ and the scale parameter $\sigma$, and a transition tensor of shape \code{(num\_actions, num\_states, num\_states)}. Every estimator returns an \code{EstimationSummary} object whose contract is described in Section~\ref{sec:design}.

\subsection{Structural likelihood} \label{sec:methods:structural}

The structural likelihood family directly maximizes the log-likelihood in equation~\eqref{eq:loglik} under the choice probabilities in equation~\eqref{eq:pi}. The challenge is that the choice probabilities depend on $V$, which depends on $r$, which depends on $\theta$. The six estimators in this family differ in how they handle this nesting.

\subsubsection{Nested fixed point (NFXP)}

\paragraph{Objective.} NFXP \citep{rust1987} maximizes the log-likelihood with the Bellman equation enforced exactly at every candidate $\theta$,
\begin{equation}
\hat\theta_{\mathrm{NFXP}} = \argmax_{\theta} \ell(\theta) \quad \text{subject to } V_\theta = T_\theta V_\theta,
\label{eq:nfxp}
\end{equation}
where $T_\theta$ is the soft Bellman operator implied by equations~\eqref{eq:v} and \eqref{eq:V} at parameter $\theta$.

\begin{algorithm}[H]
\SetAlgoLined
\KwIn{panel data, transitions $P$, features $\phi$, initial $\theta_0$}
\KwOut{$\hat\theta$, Hessian $H$}
\For{outer $k = 0, 1, \dots$ until $\theta$ converges}{
  inner solve: iterate $V \leftarrow T_{\theta_k} V$ by successive approximation\;
  switch to Newton-Kantorovich once contraction slows \citep{iskhakov2016}\;
  evaluate $\ell(\theta_k)$ and per-observation scores $\nabla_\theta \log \pi(a_{it} \mid s_{it}; \theta_k)$\;
  update $\theta_{k+1}$ by BHHH using the score outer product, or by L-BFGS-B\;
}
return $\hat\theta = \theta_k$, $H = \nabla^2 \ell(\hat\theta)$ via JAX\;
\caption{NFXP polyalgorithm}
\end{algorithm}

\paragraph{Versions.} The \code{NFXP} class exposes three inner-solve modes, namely successive approximation, Newton-Kantorovich, and the SA-to-NK polyalgorithm of \citet{iskhakov2016}. The outer optimizer defaults to BHHH with analytical per-observation scores and switches to L-BFGS-B when the user disables score caching. Standard errors come from the BHHH outer-product or from the Hessian.

\subsubsection{Mathematical programming with equilibrium constraints (MPEC)}

\paragraph{Objective.} MPEC \citep{su2012} treats $\theta$ and $V$ as joint optimization variables under the Bellman equation as an equality constraint,
\begin{equation}
\hat\theta_{\mathrm{MPEC}}, \hat V = \argmax_{\theta, V} \ell(\theta, V) \quad \text{subject to } V = T_\theta V.
\label{eq:mpec}
\end{equation}

\paragraph{Versions.} The \code{MPEC} class exposes a full-space SQP solver and a reduced-space variant that solves the Bellman equation analytically inside the subproblem. The reduced-space version is roughly three times slower on the Rust panel but does not require sparse linear algebra. JAX provides Jacobians for both modes.

\subsubsection{Conditional choice probability (CCP)}

\paragraph{Objective.} The CCP estimator \citep{hotz1993} replaces the inner Bellman solve with the Hotz-Miller inversion. Under multinomial logit errors, value differences are
\begin{equation}
v(s, a) - v(s, a^\ast) = \sigma \log \pi(a \mid s) - \sigma \log \pi(a^\ast \mid s),
\label{eq:hotz}
\end{equation}
so the structural likelihood can be evaluated from any consistent estimate of $\pi$ without solving the Bellman equation. \citet{aguirregabiria2002} pair this with nested pseudo-likelihood iterations on $\theta$ and $\hat\pi$.

\paragraph{Versions.} The \code{CCP} class exposes one-step and NPL modes. One-step CCP uses empirical choice frequencies and stops at $k = 1$. NPL iterates to joint convergence and is the default for sparse panels where empirical CCPs are noisy. Both modes self-initialize from the data, which is the practical edge over NFXP on the Rust panel.

\subsubsection{Sieve efficient estimation (SEES)}

\paragraph{Objective.} SEES \citep{luo2024sees} replaces the value function with a sieve approximation $V_\xi(s) = \sum_{k=1}^{K} \xi_k b_k(s)$ in a basis of $K$ polynomials or splines, and solves
\begin{equation}
\hat\theta_{\mathrm{SEES}}, \hat\xi = \argmax_{\theta, \xi} \ell(\theta, \xi) - \kappa \sum_s \big( V_\xi(s) - (T_\theta V_\xi)(s) \big)^2,
\label{eq:sees}
\end{equation}
where $\kappa$ is a Bellman residual penalty tuned by cross-validation.

\subsubsection{Neural network efficient estimation (NNES)}

\paragraph{Objective.} NNES \citep{nguyen2025nnes} replaces the value function with a feedforward neural network $V_\eta(s)$ and solves
\begin{equation}
\hat\theta_{\mathrm{NNES}}, \hat\eta = \argmax_{\theta, \eta} \ell(\theta, \eta) - \kappa \sum_s \big( V_\eta(s) - (T_\theta V_\eta)(s) \big)^2.
\label{eq:nnes}
\end{equation}
The estimator is root-$n$ consistent and semiparametrically efficient, with a block-diagonal information matrix between $\theta$ and $\eta$.

\subsubsection{Temporal-difference CCP (TD-CCP)}

\paragraph{Objective.} TD-CCP \citep{adusumilli2025tdccp} combines the Hotz-Miller inversion with neural function approximation and cross-fitting. The estimator splits the panel into $K$ folds by individual, fits the conditional choice probabilities $\hat\pi^{(-k)}$ on the complement of fold $k$, and uses the held-out CCPs to evaluate the score on fold $k$. The score has the form
\begin{equation}
\psi(s, a; \theta, \hat\pi^{(-k)}) = \nabla_\theta \log \pi(a \mid s; \theta) + \beta \nabla_\theta \E_{s' \sim P} \big[ \log \sum_{a'} \exp(Q(s', a'; \theta, \hat\pi^{(-k)})/\sigma) \big].
\label{eq:tdccp}
\end{equation}
Cross-fitting renders the influence function of $\hat\theta$ orthogonal to the first-stage CCP estimation error and yields valid inference under high-dimensional state spaces.

\begin{algorithm}[H]
\SetAlgoLined
\KwIn{panel, $P$, $\phi$, fold count $K$, neural CCP architecture}
\KwOut{$\hat\theta$, cross-fit Hessian $\hat H$}
split individuals into $K$ folds at random with seed 42\;
\For{$k = 1, \dots, K$}{
  fit $\hat\pi^{(-k)}$ by softmax neural classification on the complement of fold $k$\;
}
solve $\sum_{i \in k(i)} \psi(s_{it}, a_{it}; \theta, \hat\pi^{(-k(i))}) = 0$ jointly across folds\;
estimate $\hat H$ from the within-fold score outer products\;
\caption{Cross-fit TD-CCP}
\end{algorithm}

\paragraph{Versions.} The \code{TDCCP} class exposes two split rules. The default splits by individual and produces a valid cross-fit influence function. The row split is provided for ablations and breaks the orthogonality argument.

\subsection{Maximum entropy inverse reinforcement learning}

The maximum entropy inverse reinforcement learning family treats the reward as unknown and recovers it by matching the feature expectations of the expert under a maximum entropy policy.

\subsubsection{Maximum causal entropy IRL (MCE-IRL)}

\paragraph{Objective.} MCE-IRL \citep{ziebart2010} parametrizes the reward as $r(s, a) = \theta^\top \phi(s, a)$ and finds the $\theta$ that matches feature expectations under the maximum causal entropy policy,
\begin{equation}
\E_{\pi_\theta}[\phi(s, a)] = \hat \E[\phi(s, a)],
\label{eq:mce}
\end{equation}
where $\pi_\theta$ on the left is the soft-Bellman policy at $\theta$.

\begin{algorithm}[H]
\SetAlgoLined
\KwIn{panel, $P$, $\phi$, initial $\theta_0$}
\KwOut{$\hat\theta$, asymptotic $\hat H$}
compute $\hat \E[\phi]$ once from the panel\;
\For{$k = 0, 1, \dots$}{
  inner solve: soft value iteration to obtain $\pi_{\theta_k}$\;
  compute occupancy $\mu_{\pi_{\theta_k}}$ by direct linear solve, see \citet{gleave2022imitation}\;
  evaluate gradient $\hat \E[\phi] - \E_{\pi_{\theta_k}}[\phi]$\;
  update $\theta_{k+1}$ by L-BFGS-B step on the gradient\;
  stop when both gradient norm and $\| \mu_{\pi_{\theta_k}} - \hat\mu \|_\infty$ fall below tolerance\;
}
\caption{Maximum causal entropy gradient path}
\end{algorithm}

\paragraph{Versions.} The \code{MCEIRL} class exposes the gradient-path estimator with dual stopping criteria (gradient norm and occupancy distance) and the explicit feature-matching solver that solves equation~\eqref{eq:mce} as a root finding problem. The gradient path is the default. On panels with linear features and action-dependent reward, MCE-IRL produces numerically indistinguishable estimates from NFXP, which we verify in Section~\ref{sec:illustrations:mce}.

\subsection{Q-function based recovery}

Two estimators in \pkg{econirl} avoid an explicit reward parametrization and recover the choice-relevant primitive directly.

\subsubsection{Inverse soft Q-learning (IQ-Learn)}

\paragraph{Objective.} IQ-Learn \citep{garg2021} fits the soft Q-function directly by minimizing a chi-squared distance between the model-implied and empirical state-action occupancies, plus a Bellman residual penalty,
\begin{equation}
\hat Q = \argmin_{Q} \chi^2(\mu_{\pi_Q} \,\|\, \hat\mu) + \kappa \sum_{s, a} \big( Q(s, a) - r_Q(s, a) - \beta \E_{s'} V_Q(s') \big)^2,
\label{eq:iql}
\end{equation}
with $\pi_Q$ the soft policy at $Q$ and $r_Q$ the implied reward through the inverse Bellman operator.

\paragraph{Versions.} The \code{IQLearn} class exposes a tabular Q-function for finite state spaces and a neural Q head for larger ones. The chi-squared objective trades likelihood maximization for speed, and the implied reward is identified up to a state-dependent constant. Tabular Q does not propagate to unvisited states, which the user should check before relying on the recovered reward.

\subsubsection{GLADIUS}

\paragraph{Objective.} GLADIUS \citep{kang2025gladius} is a model-free estimator that parametrizes both the soft Q-function and the integrated value function by separate neural networks $Q_\eta$ and $V_\zeta$ and trains them jointly,
\begin{equation}
\hat\eta, \hat\zeta = \argmin_{\eta, \zeta} -\sum_{i, t} \log \pi_{Q_\eta}(a_{it} \mid s_{it}) + \kappa \sum_{i, t} \big( Q_\eta(s_{it}, a_{it}) - \beta V_\zeta(s_{it+1}) \big)^2.
\label{eq:gladius}
\end{equation}
The first term is the imitation log-likelihood through the soft policy implied by $Q_\eta$. The second term is a Bellman consistency penalty that uses observed transitions in the panel rather than a transition matrix.

\paragraph{Versions.} The \code{GLADIUS} class exposes a Q-only mode in which $V_\zeta$ is replaced by the soft maximum of $Q_\eta$, and a dual mode in which $Q_\eta$ and $V_\zeta$ are independent networks. The Bellman penalty weight $\kappa$ trades off imitation accuracy and Bellman consistency. \code{GLADIUS} is the only structural-style estimator in the library that operates without a supplied transition matrix, which makes it the right default for environments where the dynamics are not known.

\subsection{Adversarial inverse reinforcement learning}

\subsubsection{AIRL}

\paragraph{Objective.} AIRL \citep{fu2018} recovers the reward by training a discriminator to distinguish expert state-action pairs from pairs generated by a learned policy. The discriminator is parametrized as
\begin{equation}
D_\theta(s, a, s') = \frac{\exp(f_\theta(s, a, s'))}{\exp(f_\theta(s, a, s')) + \pi(a \mid s)},
\quad f_\theta(s, a, s') = r_\theta(s, a) + \beta h_\theta(s') - h_\theta(s),
\label{eq:airl}
\end{equation}
so the recovered $r_\theta$ is disentangled from the dynamics and transfers across transition structures. AIRLHet \citep{lee2026airlhet} adds an expectation-maximization loop over latent agent types.

\begin{algorithm}[H]
\SetAlgoLined
\KwIn{panel, environment simulator, reward net $r_\theta$, shaping net $h_\theta$, policy net $\pi_\omega$}
\KwOut{$\hat\theta$, $\hat\omega$}
\For{epoch $k = 0, 1, \dots$}{
  sample expert batch from the panel and policy batch from rollouts of $\pi_\omega$\;
  step $\theta$ to maximize the discriminator log-likelihood in equation~\eqref{eq:airl}\;
  step $\omega$ to maximize the policy reward $\log D_\theta - \log(1 - D_\theta)$ via PPO\;
}
\caption{Adversarial IRL}
\end{algorithm}

\paragraph{Versions.} The \code{AIRL} class exposes the base AIRL of \citet{fu2018} and the heterogeneous-agent extension \code{AIRLHet} of \citet{lee2026airlhet}. AIRLHet adds an EM loop over a discrete latent type, which the user configures via the number of types.

\subsubsection{$f$-IRL}

\paragraph{Objective.} $f$-IRL \citep{ni2021firl} generalizes AIRL by replacing the binary cross-entropy discriminator loss with an arbitrary $f$-divergence between the policy and expert occupancies,
\begin{equation}
\hat\theta = \argmin_\theta D_f \big( \mu_{\pi_\theta} \,\|\, \hat\mu \big),
\label{eq:firl}
\end{equation}
with $D_f$ chosen from the Kullback-Leibler, reverse-KL, Jensen-Shannon, or chi-squared family.

\paragraph{Versions.} The \code{FIRL} class exposes the choice of $f$-divergence as a string argument, with the Kullback-Leibler form as the default. The reverse-KL form recovers a mode-seeking policy and is preferred when the expert demonstrations are sparse.

\subsection{Behavioral cloning baseline}

\subsubsection{BC}

\paragraph{Objective.} BC fits a multinomial logit model to the marginal distribution of actions conditional on states without exploiting the dynamic structure,
\begin{equation}
\hat\theta_{\mathrm{BC}} = \argmax_\theta \sum_{i, t} \log \pi_{\mathrm{BC}}(a_{it} \mid s_{it}; \theta),
\label{eq:bc}
\end{equation}
where $\pi_{\mathrm{BC}}$ is a logit on the same features $\phi$ but with the reward and Bellman terms dropped.

\paragraph{Versions.} The \code{BC} class exposes a tabular logit and a neural variant in which the policy head is a small feedforward network. BC is the lower-bound baseline for every benchmark in the library. Any estimator that fails to beat BC on out-of-sample policy quality is not learning from the Markov decision process structure.

\subsection{Estimator taxonomy}

Table~\ref{tab:taxonomy} summarizes the twelve production estimators along the structural assumptions that determine which estimator applies to which problem. The columns mirror the \code{ProblemCapabilities} dataclass in \code{econirl.estimation.categories}, which serves as the canonical taxonomy and is regenerated whenever a new estimator is added to the library.

\begin{table}[t]
\centering
\small
\setlength{\tabcolsep}{6pt}
\renewcommand{\arraystretch}{1.10}
\begin{tabular}{l l l c c c c}
\toprule
\textbf{Estimator} & \textbf{Family} & \textbf{Reward} & \textbf{Trans.} & \textbf{Inner} & \textbf{Cont.} & \textbf{Recovers} \\
                   &                 & \textbf{form}   & \textbf{known}  & \textbf{solve} & \textbf{states} & \textbf{params} \\
\midrule
NFXP        & Structural likelihood & linear  & \cmark & \cmark & \xmark & \cmark \\
MPEC        & Structural likelihood & linear  & \cmark & \xmark & \xmark & \cmark \\
CCP         & Structural likelihood & linear  & \cmark & \xmark & \xmark & \cmark \\
\addlinespace[2pt]
SEES        & Sieve approximation   & linear  & \cmark & \xmark & \cmark & \cmark \\
NNES        & Neural approximation  & linear  & \cmark & \xmark & \cmark & \cmark \\
TD-CCP      & Neural CCP            & linear  & \cmark & \xmark & \cmark & \cmark \\
\addlinespace[2pt]
MCE-IRL     & Max causal entropy    & linear  & \cmark & \cmark & \xmark & \cmark \\
GLADIUS     & Q-network IRL         & neural  & \xmark & \xmark & \cmark & \cmark \\
IQ-Learn    & Q-function IRL        & tabular & \cmark & \xmark & \xmark & \xmark \\
AIRL        & Adversarial IRL       & linear  & \cmark & \cmark & \xmark & \xmark \\
$f$-IRL     & Distribution IRL      & tabular & \cmark & \cmark & \xmark & \xmark \\
\addlinespace[2pt]
BC          & Imitation             & none    & \xmark & \xmark & \xmark & \xmark \\
\bottomrule
\end{tabular}
\caption{Twelve production estimators along their structural assumptions. Columns are derived from the \code{ProblemCapabilities} dataclass in \code{econirl.estimation.categories}, which serves as the canonical taxonomy and is regenerated whenever a new estimator is added to the library.}
\label{tab:taxonomy}
\end{table}

The columns can be read as a decision aid. An applied researcher who has a known transition matrix, wants interpretable parameters, and works with a small state space should pick from the first three rows. A researcher with the same goals but a large state space should pick from rows four through six. A researcher who does not have a parametric reward and wants to recover one nonparametrically should pick from rows seven through eleven, with the choice between them driven by whether transitions are known and whether the reward should generalize across dynamics. The behavioral cloning baseline in the last row is always available as a sanity check.
